\documentclass[../principal.tex]{subfiles}
\graphicspath{{\subfix{../imagenes/}}}

\begin{document}


\chapter{Voltaje}

\section{Definición}

Las máquinas sonoras descritas en este libro emiten y reciben señales codificadas en forma de voltajes. Por eso reciben los nombres VCX, por la sigla en inglés de X controlado por voltaje.

\section{Unidades de medida}

El voltaje se mide en joule por coulomb (\unit{\joule\per\coulomb}), pero por conveniencia se usa el volt (\unit{\volt}), donde:

\begin{equation}
    \unit{\volt} = \unit{\joule\per\coulomb}
\end{equation}


Existen otras formas de transmitir información, por ejemplo, por corriente, que es un esfuerzo que está llevando a cabo Peter Blasser en esta década con amplificadores de Norton.

\section{Tierra o voltaje de referencia}

Los voltajes no son magnitudes aboslutas, son relativas.

Por ejemplo la ciudad de Santiago en Chile tiene una altitud media de 543 metros por sobre el nivel del mar. Esto es, si le asignamos el valor 0 a la altura del nivel del mar, desde ahí podemos medir la altura de Santiago y expresarla no en absoluto, sino que relativamente a este lugar de referencia.

Esto también ocurre con los voltajes, definimos un lugar como voltaje nulo, y recibe tradicionalmente el nombre de tierra, o en inglés "ground" o GND, y todos los otros voltajes de un circuito están medidos contra esta referencia arbitraria.

\end{document}


